\section{Numerical Diagnostics}

The first question in a low-rank approximation experiment is whether the
singular values decay quickly enough to justify truncation. Figure
\ref{fig:singular-value-decay} shows a toy comparison between fast and slow
decay.

\begin{figure}[h]
  \centering
  \includegraphics[width=0.72\linewidth]{singular-value-decay.pdf}
  \caption{Toy singular value decay curves used as a diagnostic for low-rank
  approximation.}
  \label{fig:singular-value-decay}
\end{figure}

The second question is whether the compressed representation saves enough
storage to justify the approximation. Figure \ref{fig:storage-vs-rank} compares
full tensor storage with a simple equal-rank Tucker storage count.

\begin{figure}[h]
  \centering
  \includegraphics[width=0.72\linewidth]{storage-vs-rank.pdf}
  \caption{A lightweight storage comparison for a cubic tensor and equal Tucker
  rank.}
  \label{fig:storage-vs-rank}
\end{figure}

The corresponding notebook keeps these diagnostics close to the code that
generates them, so that dissertation figures and computational checks can be
reviewed together.

